\chapter{Software Project Management Plan}
\label{chap:plan}

\section{Overview}

The project was run as an incremental, test-driven build: each phase produced a working, tested increment of the existing system rather than a rewrite. The plan below describes the process, organisation, schedule, risks and quality practices.

\section{Process model}

An iterative approach with short cycles of \emph{design, implement, test, review}. Each cycle ended with the automated test suite passing and a manual check of the changed workflow before the next cycle began.

\begin{figure}[H]
\centering
\begin{tikzpicture}[node distance=0.55cm]
  \node[stage] (p1) {Audit};
  \node[stage, right=of p1] (p2) {Design system};
  \node[stage, right=of p2] (p3) {Wizard};
  \node[stage, right=of p3] (p4) {Dashboards};
  \node[stage, below=0.9cm of p1] (p5) {Registration};
  \node[stage, right=of p5] (p6) {Voting and booths};
  \node[stage, right=of p6] (p7) {Results and audit};
  \node[stage, right=of p7] (p8) {Hosting, docs, QA};
  \foreach \a/\b in {p1/p2, p2/p3, p3/p4, p5/p6, p6/p7, p7/p8} {\draw[flow] (\a) -- (\b);}
  \draw[flow] (p4.south) -- ++(0,-0.45) -| (p5.north);
\end{tikzpicture}
\caption{Project phases}
\end{figure}

\section{Phases and deliverables}

\begin{center}
\small
\begin{tabularx}{\linewidth}{@{}>{\raggedright}p{3.1cm} X >{\raggedright\arraybackslash}p{3.6cm}@{}}
\toprule
\textbf{Phase} & \textbf{Activities} & \textbf{Deliverable} \\
\midrule
1. Codebase audit & Inspect architecture, authentication, schema and voting flow; list reusable parts and risks & Audit notes; list of defects found \\
2. Design system & Colours, typography, components, light/dark themes, animation rules & Reusable UI components \\
3. Installation wizard & First-run detection, validation, atomic installation, locking & Wizard and setup API \\
4. Role dashboards & Election Officer, Desk Officer and voter experiences; navigation and permissions & Three role-specific layouts \\
5. Registration & Programmes, slip ID generation, search, printing & Desk workflow \\
6. Voting & Booth flow, animations, PINs, duplicate prevention & Voting experience \\
7. Booths and results & Booth tokens, revocation, hidden results, reveal, public link, draws, audit log & Booth and results features \\
8. Hosting, documentation, QA & Docker Compose, migrations, guides, test suite, security review & Release candidate and documents \\
\bottomrule
\end{tabularx}
\end{center}

\section{Work breakdown}

\begin{itemize}
  \item \textbf{Backend:} models and migrations; setup, auth, admin, desk, booth, voter, results blueprints; vote service; audit.
  \item \textbf{Frontend:} design tokens and components; wizard; officer, desk and booth screens; results and projector view; profile.
  \item \textbf{Operations:} Dockerfiles, Compose stack, kiosk lock-down scripts, backup and upgrade guides.
  \item \textbf{Quality:} unit, integration and security tests; manual end-to-end checks; documentation.
\end{itemize}

\section{Roles and responsibilities}

\begin{center}
\begin{tabularx}{\linewidth}{@{}lX@{}}
\toprule
\textbf{Role} & \textbf{Responsibility} \\
\midrule
Developer (student) & Analysis, design, implementation, testing, documentation \\
Project guide & Review of requirements, design and deliverables \\
College election committee (stakeholder) & Provides requirements: positions, departments, procedures \\
\bottomrule
\end{tabularx}
\end{center}

\section{Tools and environment}

\begin{center}
\begin{tabularx}{\linewidth}{@{}lX@{}}
\toprule
\textbf{Purpose} & \textbf{Tool} \\
\midrule
Version control & Git and GitHub \\
Backend & Python 3.11, Flask, SQLAlchemy, pytest \\
Frontend & Node.js, React, Vite, Tailwind CSS \\
Containers & Docker and Docker Compose \\
Database & PostgreSQL (SQLite for tests) \\
Documentation & Markdown, \LaTeX{} (Overleaf), TikZ diagrams \\
\bottomrule
\end{tabularx}
\end{center}

\section{Risk management}

Risks are reviewed at the end of each phase. The main risks and their handling are listed in Chapter~\ref{chap:intro} (risk analysis). Additional project risks:

\begin{center}
\small
\begin{tabularx}{\linewidth}{@{}>{\raggedright}p{4.2cm} X@{}}
\toprule
\textbf{Risk} & \textbf{Response} \\
\midrule
Changing requirements from stakeholders & Small increments; behaviour covered by tests so changes are safe \\
Breaking existing behaviour & Existing tests kept and extended; database migrations are idempotent and preserve data \\
Environment differences (Windows, Linux) & Everything runs in Docker; a separate isolated test stack was used for manual checks \\
Time pressure & Highest-risk work (voting integrity, security) done first \\
\bottomrule
\end{tabularx}
\end{center}

\section{Quality assurance}

\begin{itemize}
  \item Automated tests run before every change is accepted (Chapter~\ref{chap:test}).
  \item Security review of authentication, authorisation, input handling and the vote path; findings were fixed and covered by tests.
  \item Manual walk-through of each role's workflow in the browser, in light and dark themes and at phone width.
  \item Documentation updated with each feature.
\end{itemize}

\section{Configuration management}

Source code is under Git. Configuration is separated from code through environment variables, with \code{.env.example} files documenting each value. Database changes are recorded as ordered, idempotent migrations (currently ten) and listed in the changelog. Releases are tagged; backups are taken before every upgrade.

\section{Effort and schedule}

\begin{center}
\small
\begin{tabularx}{\linewidth}{@{}l c X@{}}
\toprule
\textbf{Phase} & \textbf{Share of effort} & \textbf{Notes} \\
\midrule
Audit and design system & 10\% & Reuse of working parts decided here \\
Wizard and dashboards & 20\% & \\
Registration and voting & 25\% & Highest risk; most tests \\
Booths, results, audit & 20\% & Added after review of the voting flow \\
Hosting, documentation, QA & 25\% & Includes security fixes and reports \\
\bottomrule
\end{tabularx}
\end{center}

\section{Maintenance plan}

After release: keep dependencies updated; run the test suite in continuous integration on every push; follow the upgrade guide for each release; keep the changelog; publish security fixes promptly. The self-hosting guide allows any institution to maintain its own installation.
